
\section{Linear Methods in Machine Learning}
\definition{Linear Methods}{A foundational class of algorithms that model relationships between input features and output predictions by computing a weighted sum of the features. They assume the decision boundary is a straight line (2D) or a hyperplane (higher dimensions).}



\subsection{Linear Regression}
\definition{Linear Regression}{
    Linear regression predicts a continuous target variable by finding the optimal weights \f{w} and bias \f{w_0}.
    \cf{f(x_i; w) = w_0 + w_1 x_{i,1} + ... + w_D x_{i,D} = w_0 + x_i w^T}
    It assumes a strictly linear relationship between inputs and outputs, which may fail to capture complex data structures.
}
\definition{Basis Functions}{
    Mathematical transformations \f{\phi(x)} applied to input features to map them into a higher-dimensional space, allowing linear models to capture non-linear relationships (e.g., polynomial or Gaussian functions).
    \cf{f(x_i, w) = w_0 + \sum_{j=1}^M w_j \phi(x_i)_j}
}


\subsection{Optimization of Weights}
To find the best weight vector \f{w} that minimizes the loss function, two main approaches can be used:
\begin{enumerate}
\item \b{Analytical Solution}: The optimal weights can be computed directly using the normal equation (iff the errors fulfill a few assumptions). This solution becomes computationally expensive for large datasets due to the matrix inversion step, which has a complexity of \f{O(D^3)}.
\cf{w = (X^T X)^{-1} X^T y}
\item \b{Iterative Optimization (Gradient Descent)}: A numerical method used to iteratively update weights by taking steps in the direction of the steepest descent of the loss function \f{J(w)}.
\cf{w \leftarrow w - \eta \frac{\partial J(w)}{\partial w}}
\end{enumerate}



\subsection{Linear Classification \& Logistic Regression}
\definition{Linear Classification}{Categorizes data into distinct classes by defining a linear decision boundary (hyperplane) that separates the feature space.}
\definition{Logistic Regression}{
    A linear classification model used for binary classification. It applies the \b{sigmoid} (or logistic) function to squeeze the linear output into a probability between 0 and 1.
    \cf{h_w(x) = \frac{1}{1 + e^{-w^T x}}}
}



\subsection{Optimization in Logistic Regression}
Because the sigmoid function makes the loss function non-linear, logistic regression cannot be solved analytically. It must be optimized using Gradient Descent. The gradient of the binary cross-entropy loss with respect to the weights is calculated as follows:
\cf{\frac{\partial J(w)}{\partial w} = \sum_{i=1}^N (h_w(x_i) - y_i)x_i}

\definition{Stochastic Gradient Descent (SGD)}{
    An optimization variant that approximates the true gradient by using a single random data point or a small mini-batch at each step, making it much faster and more memory-efficient for large datasets.\\
    SGD introduces noise into the gradient estimation, which can prevent the model from getting stuck in local minima and allows for frequent weight updates without scanning the entire dataset.
}
